\statement{B.78}{Proved}{A complete based presentation of the actual marginal source-membership problem}
The positive marginal source-image comparison has an explicit presentation
with twenty fiber generators, two retained base generators, and all
twenty-nine filled monodromies in one based marking. The same comparison
holds for the actual closed kernel images in full completions. The
presentation below specifies the remaining uniform normality condition;
it does not establish that condition or triviality of \(\mathcal K\).

\begin{proof}
\emph{The actual surface and the two finite covers.}\par
Fix \(\zeta=\exp(2\pi i/5)\). In the positive marginal, the old source
coordinate is \(s\), and the moving fiber punctures are
\[
 p_u(s)=\left(\frac{u-s}{u+s}\right)^2\quad(u\in\mu _5),
 \qquad p=0,1,\infty.
 \tag{B.78.1}
\]
The source section is \(p=(\zeta-s)/(\zeta+s)\). Its old base excludes
\(0,\pm\mu _5\), and its eight additional deleted values satisfy
\(s^2=u(2\zeta-u)\), \(u\ne\zeta\). Set
\(q=s^2\) and \(T=s(1+p)/(1-p)\), and order the roots of
\(X^2-2TX+q\). The resulting surface is
\[
 E_{\rm ord}=\{(v,w,s):s^2=vw,\quad
 v,w\notin\{0\}\cup\mu _5,\quad v\ne w,\quad vw\notin\mu _5\}.
 \tag{B.78.2}
\]
Ordering the roots and then forgetting \(s\) are two separate unramified
double covers. The source in (B.78.2) is \(v+w=2\zeta\).
Filling the four original peripheral values indexed by \(u\ne\zeta\)
removes precisely the eight conditions \(v\ne u,w\ne u\); it retains
the diagonal, the two zero divisors, and all product hyperbolas.

Both covers extend after this filling. The ordering branch divisor
\(v=w\) and the square-root branch divisor \(vw=0\) remain deleted.
Their restrictions to the source extend for the same reason. Thus each
cover character factors through the corresponding filling map. In a
commuting square of this form, the two filling kernels are the same
subgroup of the covering fundamental group: if \(A_0\leq A\) is the
inverse image of an index-two subgroup \(B_0\leq B\) under a surjection
\(A\to B\), then
\[
       \ker(A_0\longrightarrow B_0)=\ker(A\longrightarrow B).
 \tag{B.78.3}
\]
The same statement holds in completions because finite-index subgroup
completions embed and the index-two characters are continuous. Apply
this observation to both ambient and source squares. It retains the
entire original source kernel, including conjugates from the other
ordering coset; it does not replace that kernel by one generated from
only a subset of source conjugators.

In the original marginal bundle the base is unchanged by filling.
Its free fiber groups have ranks seven and three, and its base has
vanishing \(\pi _2\). Consequently its ambient filling kernel is exactly
the fiber filling kernel. This remains true for the actual completed
kernels: every finite quotient of the finitely generated free fiber
is refined by a characteristic finite-index quotient, whose finite
automorphism action extends to a finite quotient of the split total
group. This proves the full induced fiber topology directly. Therefore
(B.78.3) identifies the actual marginal source-kernel comparison with
the one in the following polynomial surface, discretely and in full
completions. No topology on the abstract infinite-rank deletion kernel
has been assumed.

\emph{The polynomial pencil and its complete critical set.}\par
Scale \(v,w,s\) by \(\zeta\), and put
\[
 \lambda=\frac{v-1}{w-1},\qquad t=1+\lambda^{-1},\qquad z=w-1.
 \tag{B.78.4}
\]
The base is \(B=\mathbb C\setminus\{1,2\}\), the source is \(t=0\),
and the twenty generic finite punctures of its \(z\)-fiber are
\[
\begin{split}
 &0,-1,1-t,-t;\\
 &u-1,\ (u-1)(t-1)\quad(u\in\mu _5\setminus\{1\});\\
 &z^2+tz+(t-1)(1-q)=0
                   \quad(q\in\mu _5\setminus\{1\}).
\end{split}
 \tag{B.78.5}
\]
The eight points in the second line are the deleted set \(U\).
The other twelve points are retained. At \(t=0\), only \(0,-t\)
merge, so the special source and its filling have free groups of
ranks nineteen and eleven.

There are exactly twenty-nine exceptional values in \(B\):
\[
\begin{array}{c|c|c}
\text{type}&\text{values}&\text{number}\\\hline
\text{retained node}&0&1\\
\text{retained branch}&t^2-4(1-q)t+4(1-q)=0&8\\
U\text{--zero node}&2-u,\quad1+(1-u)^{-1}&8\\
U\text{--}U\text{--product triple}&
       1+(u-1)/(v-1),\quad u\ne v&12
\end{array}
 \tag{B.78.6}
\]
Here \(u,v,q\) run through the four nonidentity fifth roots. The four
triples with \(uv=1\) are \(t=1-u\). Exact arithmetic in
\(\mathbb Q[\zeta]\) proves that all displayed values are distinct,
that the branch quadratics are separable, and that no further finite
collision occurs. Two different product hyperbolas cannot intersect
in a fiber; a product puncture cannot meet a zero puncture; and all
remaining pair collisions are enumerated by (B.78.6). The monic
quadratic equations and affine rational punctures have constant fiber
degree. Thus there are no additional finite escape values.

\emph{An explicit common base and source inclusion.}\par
Use \(t_*=i/100,z_*=100\), with projection \(Pz\), where
\(P=1+371i/1000\). To compare with the original source base \(s=2+2i\),
use the polygon
\[
        2+2i\longrightarrow10+10i
                 \longrightarrow\zeta i\sqrt{9999}.
 \tag{B.78.7}
\]
All deleted source points have modulus less than two. The first segment
lies outside that disk; the second has imaginary part at least ten,
since \(\Re(\zeta)>3/10\). At its endpoint choose
\(v=-99\zeta,w=101\zeta\) and lift the path backwards to fix the
original ordering. After scaling, continue along \(t=iy\),
\(0\leq y\leq1/100\), with \(z=100\), lifting \(s^2=vw\).
This is a path in the actual spaces and covers.

Use the standard Artin cut system in the projected coordinate, with
pullback convention
\[
 \Phi_{\sigma_i}(y_i)=y_i y_{i+1}y_i^{-1},\qquad
 \Phi_{\sigma_i}(y_{i+1})=y_i,
 \tag{B.78.8}
\]
and all other generators fixed. The based positive boundary of a disk
enclosing consecutive points \(i,\ldots,j\) is
\(y_i\cdots y_j\). This fixes both the whisker and action conventions;
the transported meridian along a whisker \(W\) is
\(\Phi_W\) of the local meridian. In particular this is a pullback
action, not the inverse convention for the last-strand point-pushing
conjugator encountered earlier.

Along the short path in (B.78.7), the disk \(D=\{|Pz|<1/20\}\)
contains precisely the merging pair. All other projected real
coordinates stay outside \((-1/100,1/100)\), and in fact their
absolute values exceed \(19/100\). Choose the common based cut to the
lower boundary of \(D\), and transport the other cuts outside \(D\).
The merging points have generic indices ten and eleven. Define the
special meridian using that same cut and the boundary of \(D\).
It is literally \(y_{10}y_{11}\), rather than merely conjugate to it.
Thus, with
\[
 x_{10}=c=y_{10}y_{11},\qquad b=x_{11}=y_{11},\qquad
 x_i=y_i\ (i\ne10),
 \tag{B.78.9}
\]
the actual based special source is represented by the free factor
\[
 H=\langle x_i:i\ne11\rangle\leq F_{20}=H*\langle b\rangle,
 \qquad \mathcal U=\{1,2,6,9,12,15,19,20\}.
 \tag{B.78.10}
\]
Here \(\mathcal U\) lists the eight deleted meridians. This is a chosen
basis of the actual source, not an identification with an independently
chosen old lasso basis. Its based image and filling kernel are unchanged
by that explicit basis choice. The local relation at zero is
\([y_{10},y_{11}]=1\), equivalently \([c,b]=1\).

\emph{The complete exact braid certificate.}\par
For each of the thirty-one values consisting of (B.78.6) and \(1,2\),
take a square of half-side \(1/1000\) about the rational center specified
in \path{work_DE/next_marginal_source/complete_pencil_words.json}.
The centers approximate the actual algebraic values with error less
than \(10^{-25}\). A lasso goes straight from \(i/100\) to its northeast
corner, goes counterclockwise around the square, and returns. Exact
distance and determinant tests prove that the squares are disjoint,
each straight whisker misses every other square, and all whisker
directions are distinct. These are consequently a distinguished system
of generators for the thirty-one-punctured base, not arbitrary
conjugates of a proposed generating system.

The initial projected strand-label order is
\[
 (6,8,15,17,13,4,1,19,10,0,3,11,18,2,5,12,16,14,9,7).
 \tag{B.78.11}
\]
Labels zero through three are the first line of (B.78.5); labels four
through eleven are its constant/moving \(U\) pairs in increasing powers
of \(\zeta\); the remaining labels are the two branches of each
quadratic. The certificate records all thirty-one words
\(\beta_d\in B_{20}\), including every whisker and local circle.
Their local exponent checks are
\[
\begin{array}{c|r|r}
\text{lasso type}&\text{number}&\text{exponent per lasso}\\\hline
t=1&1&128\\
t=2&1&26\\
t=0&1&2\\
\text{retained branch}&8&1\\
U\text{--zero node}&8&2\\
\text{triple}&12&6
\end{array}
 \tag{B.78.12}
\]
The eight branch words have the expected transpositions; all other
words are pure. Exponents are consistency checks, not the proof of
the words or of the presentation.

For completeness, the crossing method is exact. On every straight
segment a projected quadratic pair has form \(c\pm\sqrt{A+iB}\),
with rational polynomial coefficients. The candidate equations are
linear for two rational points, \(B=0\) for a quadratic self-crossing,
\(4h^2(h^2-A)-B^2=0\) for a rational/quadratic crossing, and
\[
 N^2-4AN\delta-4B^2\delta^2=0,
 \quad \delta=C-A,\quad N=B^2-D^2,
 \tag{B.78.13}
\]
for two pairs with squares \(A+iB,C+iD\). Their degrees are at most
six. Rational Sturm intervals isolate every candidate; nested rational
square-root bounds certify every order and crossing height. Odd
negative-real-axis crossings track analytic branches. Simultaneous
crossings are allowed only for disjoint strand pairs. Endpoint orders
and loop permutations are independently checked.

The cyclotomic coefficients have certified rational approximations
within \(10^{-30}\) per component. Exact polynomial lower bounds along
all segments give pair separation greater than \(10^{-14}\); coefficient
deformation moves every strand by less than \(10^{-22}\). Thus the
words are those of the actual cyclotomic family. This argument deforms
only the loop configurations and does not assume that rational
perturbation preserves the entire discriminant arrangement.
The producers and genuine outputs are
\path{work_DE/next_marginal_source/quadratic_pencil_braid.py},
\path{work_DE/next_marginal_source/all_filled_pencil_braids.py},
\path{work_DE/next_marginal_source/audit_pencil_geometry.py}, and
\path{work_DE/next_marginal_source/audit_complete_pencil.py},
with their respective \texttt{.out} and \texttt{.err} files. The two
full braid JSON files retain all 155 segment root-isolation and height
records; the combined word file includes hashes of its inputs.

\emph{The finite presentation and the precise remaining gate.}\par
Let \(\mathscr D\) be the twenty-nine filled values, and let
\(\beta_1,\beta_2\) denote the two retained lassos. The actual total
group is
\[
\begin{split}
G=\langle y_1,\ldots,y_{20},\tau_1,\tau_2\mid\;&
 \tau_jy_i\tau_j^{-1}=\Phi_{\beta_j}(y_i)
                   \quad(j=1,2;\ 1\leq i\leq20),\\
 &\Phi_{\beta_d}(y_i)=y_i
                   \quad(d\in\mathscr D;\ 1\leq i\leq20)\rangle.
\end{split}
 \tag{B.78.14}
\]
The stable letters use the positive section loops with the pullback
convention (B.78.8). Here is the geometric sign check. If \(h_s\) is
the puncture isotopy along a positive base loop, moving a based fiber
loop \(\gamma\) through its swept square gives the boundary relation
\(\gamma\tau h_1(\gamma)^{-1}\tau^{-1}=1\). Therefore
\(\tau^{-1}\gamma\tau=(h_1)_*(\gamma)\), or equivalently
\(\tau\gamma\tau^{-1}=(h_1^{-1})_*(\gamma)=\Phi(\gamma)\).
For a positive local half twist the left strand moves through the lower
half-plane; the standard cuts in (B.78.8) give precisely this inverse
transport. Pullbacks compose in the recorded traversal order, which
also explains why a whisker uses \(\Phi_W\) on its local meridians.
One may equivalently reverse both the stable-letter and monodromy
conventions; they must not be changed independently.

Indeed, over the punctured base the group is the split bundle group.
A far-point section exists on the entire base and may be chosen to
equal \(z=100\) on all the certified loops: for example take
\(z=100\max(1,|t|/10)\), which exceeds every puncture modulus.
Filling each of the twenty-nine parameter disks kills its section
loop and imposes exactly its variation relations. Van Kampen gluing
therefore gives (B.78.14). The generic-fiber comparison also satisfies
\citet[Corollary~1.1 and Remark~1.2]{DE_Shimada2010}: the open surface
projection is a submersion, every fiber is a nonempty irreducible
punctured affine line, and the base has vanishing \(\pi_2\).
No special-fiber surjectivity is inferred from that theorem.

Put
\[
 N=\langle\!\langle x_i:i\in\mathcal U\rangle\!\rangle_H,
 \quad I=\operatorname{im}(N\to G),
 \quad R=\langle\!\langle x_i:i\in\mathcal U\rangle\!\rangle_G.
 \tag{B.78.15}
\]
Filling the eight divisors gives \(G_U=G/R\): meridians normally
generate the deletion kernel by general-position disks, also at the
finitely many divisor intersections. Thus the actual marginal equality
is \(I=R\). In full completions the exact comparison is
\[
 \mathcal I=\overline I\quad\subseteq\quad
 \mathcal R=\ker(\widehat G\longrightarrow\widehat{G_U})
             =\overline R.
 \tag{B.78.16}
\]
Here \(\mathcal I\) is the compact image of the actual source
completion kernel. It is not the full completion of the abstract
group \(I\). Equality in (B.78.16) is equivalent to
\[
       q(I)=q(R)\quad\text{for every finite quotient }q:G\to Q.
 \tag{B.78.17}
\]
This follows by separation of closed subsets in a profinite group.

There is also an exact single-normalizer formulation. Discretely,
\(I=R\) if and only if \(bIb^{-1}=I\). To prove the nontrivial
direction, \(H\) already normalizes its kernel image; adding \(b\)
generates the generic fiber. Hence \(b\)-normality makes \(I\) equal
to the image of the generic-fiber normal closure of all eight deleted
meridians. Every actual forgetting monodromy preserves this normal
closure, so the two retained stable letters preserve its image as well.
Thus \(I\) is normal in \(G\) and equals \(R\). The same argument
with closures proves
\[
             \mathcal I=\mathcal R
       \quad\Longleftrightarrow\quad b\mathcal I b^{-1}=\mathcal I.
 \tag{B.78.18}
\]
Closedness of the normalizer supplies the entire procyclic closure of
\(b\); no passage from discrete injectivity to completed injectivity
is used. B.77 shows explicitly why normalization of only the four
designated peripheral pairs is insufficient. The full uniform
condition (B.78.18), and then the original joint-source comparison,
remain unresolved.
\end{proof}
